\chapter{The APDU channel}
\label{ch:apdu}

\chapterbrief
  {Send your first real DESFire command and get an answer.}
  {Composition and layering --- the last rung of the ladder.}
  {A real DESFire reporting its own chip type.}

\section{The problem}

The card is activated. Now you want to send it commands. Those are
\textbf{APDUs} (Application Protocol Data Units), and they need wrapping:

\begin{center}
\begin{tikzpicture}[font=\footnotesize\sffamily, node distance=6mm,
   r/.style={draw=brewergrey, rounded corners=2pt, inner sep=4pt, minimum height=7mm}]
  \node[r] (a) {\code{90 60 00 00 00}};
  \node[font=\scriptsize\itshape, color=brewergrey, right=3mm of a] {your APDU};
  \node[r, below=7mm of a] (b) {\code{02} \enspace \code{90 60 00 00 00}};
  \node[font=\scriptsize\itshape, color=brewergrey, right=3mm of b]
       {wrapped in an I-block (Ch.~\ref{ch:blocks})};
  \node[r, below=7mm of b] (c) {\code{02} \enspace \code{90 60 00 00 00} \enspace \code{crc crc}};
  \node[font=\scriptsize\itshape, color=brewergrey, right=3mm of c]
       {CRC appended (Ch.~\ref{ch:reader})};
  \draw[-{Stealth[length=4pt]}, brewernavy] (a) -- (b);
  \draw[-{Stealth[length=4pt]}, brewernavy] (b) -- (c);
\end{tikzpicture}
\end{center}

Someone has to own that wrapping --- and the block-number bookkeeping that goes
with it.

\subsection{The block number: a one-bit conversation counter}

ISO-DEP puts a single bit in the I-block PCB that \textbf{alternates
0,\,1,\,0,\,1\ldots} with each exchange. It exists to disambiguate
retransmissions: if a reply is lost and you resend, the card can tell ``this is
the same command again'' from ``this is the next command''.

You must toggle it after every exchange. Get out of sync and the card ignores
you --- which looks exactly like a dead card.

\begin{concept}{State that should be impossible to get wrong}
If the block number were a parameter every caller had to pass, someone would
eventually pass the wrong one.

Hide it inside the channel as \code{private} state and toggle it automatically.
The bug becomes \emph{unrepresentable} rather than merely unlikely. That is the
strongest form of encapsulation: not ``callers should not do this'' but
``callers cannot''.
\end{concept}

\subsection{DESFire's APDU wrapping}

DESFire natively speaks its own short frames, but also accepts them wrapped in
ISO-7816 APDUs, which is what we use:

\begin{center}
\begin{tikzpicture}[font=\scriptsize\sffamily, x=13mm, y=8mm]
  \foreach \i/\lbl/\c in {0/{\code{90}}/brewerblue!12, 1/{\code{INS}}/breweramber!14,
                          2/{\code{00}}/black!5, 3/{\code{00}}/black!5,
                          4/{\code{Lc}}/brewergreen!12, 5/{\code{data...}}/brewergreen!8,
                          6/{\code{00}}/black!5}{
     \fill[\c] (\i,0) rectangle ++(1,1);
     \draw[brewergrey] (\i,0) rectangle ++(1,1);
     \node at (\i+0.5,0.5) {\lbl};
  }
  \node[color=brewergrey, font=\tiny, anchor=north] at (0.5,-0.15) {CLA};
  \node[color=brewergrey, font=\tiny, anchor=north] at (1.5,-0.15) {command};
  \node[color=brewergrey, font=\tiny, anchor=north] at (3.0,-0.15) {P1, P2};
  \node[color=brewergrey, font=\tiny, anchor=north] at (4.5,-0.15) {length};
  \node[color=brewergrey, font=\tiny, anchor=north] at (6.5,-0.15) {Le};
\end{tikzpicture}
\end{center}

\hex{0x90} means ``DESFire native command, wrapped''. Replies end in a two-byte
status word: \hex{91 00} for OK, \hex{91 AF} for ``more frames follow'', and
various \hex{91 xx} errors. Appendix~\ref{app:vectors} lists them all.

\section{Build it}

\subsection{\texttt{ApduChannel} --- another seam}

\begin{jcode}
package com.midnightbrewer.hardware;

/**
 * "Something that can carry an APDU to a card and bring back the reply."
 *
 * The same trick as SpiLink, one layer up. The DESFire logic you write in
 * Part IV will depend only on THIS -- so it can be tested against a fake
 * channel replaying canned card replies, with no reader and no card.
 *
 * That is not a nicety. Crypto is the one thing you cannot debug by staring
 * at it, and this interface is what makes it testable offline.
 */
public interface ApduChannel {

    /** Send one APDU, return the response (including its status word). */
    byte[] transceive(byte[] apdu) throws SpiException;
}
\end{jcode}

One method. That is the whole contract, and it is enough to make all of Part~IV
testable.

\subsection{\texttt{InformationBlock.encode} --- the block gains behaviour}

Back in Chapter~\ref{ch:blocks} an I-block could only \code{describe()} itself.
Now it does real work:

\begin{jcode}
    /**
     * Encode this I-block into a frame: its PCB, then the data it carries.
     */
    public byte[] encode(byte[] inf) {
        // TODO: byte[] frame of length inf.length + 1;
        //       frame[0] = (byte) pcb();
        //       copy inf in after it; return frame.
        return null;
    }
\end{jcode}

\begin{deeper}{Designs pay off later than you would like}
The class you built two chapters ago, which looked like an academic exercise in
polymorphism, is now the thing that formats your wire frames.

That is the normal shape of good design: the structure was right before you knew
what you would need it for. It is also why ``you aren't gonna need it'' is
advice to apply to \emph{features}, not to \emph{structure}.
\end{deeper}

\subsection{\texttt{IsoDepChannel}}

\begin{jcode}
package com.midnightbrewer.hardware;

import java.util.Arrays;

/**
 * The ISO-14443-4 (ISO-DEP) channel: one APDU at a time, wrapped in an
 * I-block.
 *
 * The OOP on show:
 *   - COMPOSITION: it HAS-A reader (a channel is not a kind of reader).
 *   - ENCAPSULATED STATE: the block number toggles and nobody outside can
 *     touch it, or forget to.
 *   - it USES the Chapter 4 hierarchy: an InformationBlock encodes the frame.
 *
 * Simplifications, fine for DESFire's small frames: one I-block per APDU, no
 * chaining of large payloads, and no S-block "please wait" handling.
 */
public class IsoDepChannel implements ApduChannel {

    /** I-block PCB base; the low bit carries the block number. */
    private static final int IBLOCK_PCB = 0x02;

    private final Iso14443Reader reader;  // composition
    private int blockNumber = 0;          // private mutable state, owned here

    public IsoDepChannel(Iso14443Reader reader) {
        this.reader = reader;
    }

    /**
     *   1. wrap the APDU in an I-block for the current block number
     *   2. send it via the reader (which appends CRC and transceives)
     *   3. toggle the block number for next time
     *   4. the reply is [PCB, ...response...] -- strip the PCB
     */
    @Override
    public byte[] transceive(byte[] apdu) throws SpiException {
        // TODO 1: InformationBlock out =
        //             new InformationBlock(IBLOCK_PCB | blockNumber);
        //         byte[] frame = out.encode(apdu);
        // TODO 2: byte[] reply = reader.exchange(frame);
        // TODO 3: blockNumber ^= 1;    // toggle -- XOR with 1 flips a bit
        // TODO 4: if (reply.length < 1)
        //             throw new SpiException("empty ISO-DEP reply");
        //         return Arrays.copyOfRange(reply, 1, reply.length);
        return new byte[0];
    }
}
\end{jcode}

\subsection{\texttt{Iso14443Reader.exchange} --- the reusable primitive}

You have written ``append a CRC and transceive'' twice already, inside
\code{select()} and \code{rats()}. Pull it out:

\begin{jcode}
    // =================================================================
    // exchange() -- send a payload as a CRC-protected T=CL frame, return the
    // reply with its CRC stripped.
    //
    // The ISO-DEP channel sits on top of this: it hands us [PCB + apdu], we
    // handle the CRC and the wire.
    // =================================================================
    public byte[] exchange(byte[] payload) throws SpiException {
        // TODO 1: byte[] crc = calculateCrc(payload);
        // TODO 2: build frame = payload followed by crc[0], crc[1].
        // TODO 3: byte[] reply = transceive(frame, 0);
        // TODO 4: if (reply.length < 2) return new byte[0];
        //         return Arrays.copyOf(reply, reply.length - 2);
        return new byte[0];
    }
\end{jcode}

\section{The layer cake, complete}

You have now built six layers, each speaking only to its neighbours.

\begin{center}
\begin{tikzpicture}[font=\footnotesize\sffamily,
  L/.style={draw=brewernavy, rounded corners=2pt, minimum width=54mm,
            minimum height=8mm, inner sep=3pt},
  S/.style={draw=brewerblue, fill=brewerblue!8, rounded corners=2pt,
            minimum width=54mm, minimum height=7mm, inner sep=3pt},
  note/.style={font=\scriptsize\itshape, color=brewergrey, anchor=west}]
  \node[L] (df) {\code{DesfireCard}};
  \node[note, right=3mm of df] {``debit 1.50 from the wallet''};
  \node[S, below=2.4mm of df] (ac) {\code{ApduChannel}};
  \node[note, right=3mm of ac] {seam: makes crypto testable};
  \node[L, below=2.4mm of ac] (idc) {\code{IsoDepChannel}};
  \node[note, right=3mm of idc] {``wrap this APDU in an I-block''};
  \node[L, below=2.4mm of idc] (ir) {\code{Iso14443Reader}};
  \node[note, right=3mm of ir] {``append a CRC, send this frame''};
  \node[L, below=2.4mm of ir] (rr) {\code{Rc522Reader}};
  \node[note, right=3mm of rr] {``load the FIFO, press send''};
  \node[L, below=2.4mm of rr] (rd) {\code{Rc522Driver}};
  \node[note, right=3mm of rd] {``write register \hex{0x0D} = \hex{0x80}''};
  \node[S, below=2.4mm of rd] (sl) {\code{SpiLink}};
  \node[note, right=3mm of sl] {seam: makes the driver testable};
  \node[L, below=2.4mm of sl] (p4) {\code{Pi4jSpiLink}};
  \node[note, right=3mm of p4] {``shift these two bytes''};
  \foreach \a/\b in {df/ac, ac/idc, idc/ir, ir/rr, rr/rd, rd/sl, sl/p4}{
     \draw[-{Stealth[length=3.5pt]}, brewergrey] (\a) -- (\b);}
\end{tikzpicture}
\end{center}

\textbf{Every arrow points down, and there are exactly two seams.} The value of
that shape: any layer can be tested by faking the layer below, and any layer can
be replaced without touching the others. Swap \code{Pi4jSpiLink} for an FT232H
and nothing above changes. Swap the RC522 for a PN532 and you write one new
\code{Iso14443Reader} subclass; nothing above changes.

\begin{concept}{What ``good architecture'' actually buys}
It is not a matter of taste. It is the property that lets you \textbf{debug one
layer at a time instead of debugging six at once.}

When something breaks in Part~IV, you will walk down this diagram asking each
layer ``are you still telling the truth?'' --- and the first \emph{no} is your
bug. Without the seams, every failure is a whole-system failure.
\end{concept}

\section{Verify --- your first conversation with a DESFire}

\begin{jcode}
IsoDepChannel channel = new IsoDepChannel(reader);
// ISO-7816-wrapped DESFire GetVersion: 90 60 00 00 00
byte[] getVersion = { (byte) 0x90, 0x60, 0x00, 0x00, 0x00 };
byte[] resp = channel.transceive(getVersion);
System.out.println("GetVersion resp = " + hex(resp));
\end{jcode}

\begin{verifybox}
\begin{trace}
GetVersion resp = 04 01 01 33 00 18 05 91 AF
\end{trace}
\end{verifybox}

Decode it --- this is your card's identity document:

\begin{center}\small
\begin{tabular}{@{}lll@{}}
\toprule
\thead{Byte} & \thead{Value} & \thead{Meaning} \\
\midrule
\hex{04}     & vendor       & NXP \\
\hex{01 01}  & type/subtype & DESFire \\
\textbf{\hex{33}} & \textbf{version} & \textbf{major 3 = EV3} (\hex{0x12} = EV1, \hex{0x22} = EV2) \\
\hex{00 18}  & storage      & \hex{0x18} $\rightarrow$ 8\,KB \\
\hex{05}     & protocol     & ISO-14443-2/3/4 \\
\hex{91 AF}  & status       & \textbf{additional frame available} \\
\bottomrule
\end{tabular}
\end{center}

\begin{trap}{\hex{91 AF} is not an error}
It is the card saying ``I have more to tell you''. \code{GetVersion} chains three
frames --- hardware info, software info, then UID and batch number. You would
send \hex{90 AF 00 00 00} twice more to drain them.

That matters practically: if you leave the card mid-chain and send a
\emph{different} command, some cards abort with an error. Either drain the chain
or power-cycle the field before the next command. This will bite you in
Chapter~\ref{ch:auth} if you do not.
\end{trap}

If the \hex{33} byte reads \hex{0x12} or \hex{0x22} you have an EV1 or EV2. EV2
works with everything in this book. \textbf{EV1 does not} --- it has no
\code{AuthenticateEV2First}, and needs the legacy \hex{0x0A}/\hex{0xAA}
authentication, which is a different algorithm.

\section{Traps}

\begin{traptable}
Empty reply & Block number out of sync, or RATS never ran. \\
Works exactly once per power-up & Not toggling the block number. \\
\hex{91 AE} on a later command & Card was mid-chain from \hex{91 AF}. \\
\hex{6E 00} / \hex{6D 00} instead of \hex{91 xx} & ISO-7816 error --- check the CLA \hex{90} and the trailing \hex{00}. \\
Reply is one byte too long & Forgot to strip the PCB. \\
\end{traptable}

\section{Commit}

\begin{shell}
git add -A
git commit -m "M6: ApduChannel seam + IsoDepChannel; GetVersion works on silicon"
\end{shell}
